% ================================= DISCUSSAO =================================
\section{Discussão}\label{sec:discussao}

O resultado principal é negativo e consistente: nenhum modelo de AM demonstrou ganho de discriminação sobre uma regressão logística com as cinco variáveis do Escore de Mayo, e as diferenças em relação à referência, entre $-$0,013 e +0,009 de AUROC, são menores que a variação entre partições (desvios de 0,035 a 0,051). Esse resultado responde à pergunta do título num sentido prático. Com as mesmas cinco variáveis, o AM não exige dados extras, mas troca uma fórmula transparente por um modelo opaco sem ganho demonstrado. Passar para 17 variáveis acrescenta, a qualquer família, o custo de coleta e a menor elegibilidade sem imputação (de 310 para 210 pacientes com dados completos), também sem ganho demonstrado. O fato de a regressão penalizada com 17 variáveis ter tido a maior AUROC nominal indica que, nestes dados, as árvores não extraíram das variáveis extras mais informação do que um modelo linear; o desenho não testa se as duas famílias as aproveitam igualmente.

O achado está alinhado com a revisão de \citeonline{christodoulou2019systematic}, que não encontrou vantagem do AM sobre a regressão logística nas comparações com baixo risco de viés. Há uma explicação plausível para isso neste problema: com 312 pacientes e 115 eventos, o espaço para aprender interações e não linearidades que um modelo linear em escala logarítmica não capture é pequeno, e o próprio Escore de Mayo já incorpora as transformações que tornam a relação aproximadamente linear. Essa explicação, no entanto, é uma hipótese; o desenho não a testa. Os números relatados nos trabalhos que seguem o enquadramento do Kaggle, com acurácias próximas de 90\% \cite{shakir2025cirrhosis,tu2025cirrhosis}, não são comparáveis com os deste estudo, porque usam outro desfecho (três classes ou óbito sem horizonte), outros dados (sintéticos) e outro protocolo de avaliação.

A calibração diferenciou os modelos mais do que a discriminação. A probabilidade original do Escore de Mayo teve AUROC média de 0,910, obtida com coeficientes estimados nos próprios pacientes de teste, mas subestimou o risco (CITL de +0,33), o que se traduziu em 46 óbitos não detectados no limiar de 0,5, contra 26 da regressão reestimada. Uma explicação possível, não testada, é a definição do rótulo: ao excluir os pacientes censurados antes de cinco anos, a amostra rotulada fica com mais óbitos do que a coorte de onde a sobrevida basal foi estimada. Recalibrado no treino, o escore teve estimativas de calibração próximas das ideais (CITL de 0,00 e inclinação de 1,10). Entre os modelos treinados, a regressão penalizada e o \textit{random forest} com 17 variáveis produziram probabilidades comprimidas demais (inclinações de 1,74 e 1,57), um problema que a AUROC não revela e que afetaria o uso clínico das probabilidades \cite{vancalster2019calibration}. Na validação quase externa, a regressão penalizada escolheu uma penalização forte, e sua inclinação de calibração chegou a 2,27.

O braço de controle mostra o tamanho do problema do vazamento neste conjunto. Incluir o tempo de seguimento como entrada reduziu o \textit{log loss} em 14\% e elevou a AUROC macro em 0,034, um ganho maior do que qualquer diferença entre algoritmos observada na análise principal. Um estudo que compare algoritmos nesse enquadramento pode, portanto, atribuir ao algoritmo um desempenho que vem da informação vazada \cite{kapoor2023leakage}.

O estudo tem limitações importantes. Os dados vêm de um único centro, nas décadas de 1970 e 1980, antes do uso do ácido ursodesoxicólico, e não representam os pacientes atuais. Não houve coorte externa verdadeira: a validação quase externa usa pacientes do mesmo centro e período, com apenas 68 pacientes, e os coeficientes publicados de Mayo foram estimados nesses dados, o que limita o escore publicado ao papel de referência histórica. A classificação com horizonte fixo exclui 106 pacientes censurados ou transplantados antes de cinco anos; a análise de sobrevivência com riscos competitivos é a extensão natural. Os intervalos de confiança são condicionais às predições fora da partição e não incorporam a variação de refazer o treino, o que os torna exploratórios, e não houve teste de equivalência. Por fim, o orçamento de busca de hiperparâmetros foi modesto (12 configurações), e a conclusão vale para os algoritmos, as grades e o orçamento avaliados. Estes resultados descrevem desempenho retrospectivo e não constituem evidência de utilidade clínica.
